\chapter{Las preguntas incómodas, primero}
\label{cap:s-incomodas}
Antes de explicar los modelos, conviene responder las preguntas que un profesor hará al ver la carpeta del proyecto. Son
las que más preocupan, y tienen respuestas claras.

\section{¿Por qué aparece ``Claude'' en el proyecto?}
Porque el proyecto se programó con la ayuda de \textbf{Claude Code}, un asistente de programación con inteligencia
artificial. Por eso:
\begin{itemize}
\item cada cambio guardado en el historial (los \emph{commits} de git) dice ``Co-Authored-By: Claude'';
\item existe un archivo \codigo{CLAUDE.md}: son las \textbf{reglas que Brandon le puso al asistente} (no inventar datos,
  ir fase por fase, explicar cada decisión con una cifra real y que Brandon decide).
\end{itemize}

\textbf{Cómo se repartió el trabajo.} El asistente escribió la mayor parte del código y de los documentos. Brandon puso las
reglas, revisó los resultados y \textbf{tomó todas las decisiones de fondo}. El método fue siempre el mismo:
\begin{enumerate}
\item el asistente explicaba qué había que decidir, con la cifra real que lo motivaba;
\item presentaba 2 o 3 opciones con sus pros y contras;
\item \textbf{Brandon elegía};
\item la decisión quedaba escrita en una nota de \codigo{docs/decisiones/}, con sus palabras, lo que se descartó y por qué.
\end{enumerate}
Hay 25 notas de decisión y más de 40 decisiones registradas. Algunos ejemplos de lo que Brandon decidió:
\begin{itemize}
\item enfocarse en solo 5 lugares y sacar todo lo que tuviera sargazo o relación con Tulum;
\item que los ceros imposibles se traten como dato faltante y no como cero;
\item el criterio para elegir el modelo del Radar: ``el que más acierta y más cambios anticipa'';
\item repartir el dinero ``proporcional al espacio libre'';
\item el corte de ``clima raro'' cuando el método automático marcaba casi la mitad de las semanas;
\item el nombre y el tono de la campaña: ``combina las 3, porque es un orgullo ser méxicano''.
\end{itemize}
Y cuando el asistente se equivocó, Brandon lo detectó: por ejemplo, señaló que las fotos de comida no eran de Quintana Roo
(eran de Mérida y Campeche) y se retiraron.

\begin{teprofe}
\emph{``¿Esto lo hiciste tú o la IA?''} Respuesta honesta y recomendada: \emph{``Usé un asistente de programación con IA
como herramienta, igual que se usa una calculadora o una librería. Las reglas, las decisiones y la revisión son mías y
están documentadas con fecha; puedo explicar cada parte y por qué se eligió.''} Conviene decirlo antes de que lo
pregunten, y revisar qué dice el profesor o la UNRC sobre el uso de IA.
\end{teprofe}

\section{¿Por qué hay tantos archivos de código?}
Porque cada archivo hace \textbf{una sola cosa}, como los cajones de un taller: en uno están los desarmadores, en otro las
llaves. Si algo falla, se sabe en qué cajón buscar, y cada pieza se puede probar sola.

Son \textbf{56 archivos} de Python en 7 cajones (carpetas) dentro de \codigo{backend/torre/}:

\begin{center}
\small
\begin{tabularx}{\linewidth}{lrY}
\encabezado \h{Cajón} & \h{Archivos} & \h{Qué guarda, con ejemplos} \\
\codigo{base} & 23 & Descargar los datos oficiales (\codigo{ingesta_datatur.py}), limpiarlos uno por fuente
(\codigo{silver_inah.py}, \codigo{silver_denue.py}\dots: 12 archivos de limpieza, porque hay 12 fuentes con reglas
distintas) y armar el almacén. \\
\filagris\codigo{radar} & 7 & Medir dónde hay presión: el índice, la predicción, Markov, el agrupamiento. \\
\codigo{pronostico} & 7 & La forma del año, los 5 modelos, el rango del 90\,\%, las tormentas y los escenarios. \\
\filagris\codigo{campana} & 9 & El reparto del presupuesto, las reseñas, la marca, los lugares para recomendar, las fotos. \\
\codigo{envivo} & 3 & La Torre: las señales de cada semana, las reglas y la reproducción con Spark. \\
\filagris\codigo{api} & 2 & Lo que alimenta la página y el servidor que calcula en vivo. \\
\codigo{documento} & 5 & Los PDF, las gráficas, las capturas de la página, la revisión de accesibilidad y la carpeta \codigo{Entregable}. \\
\bottomrule
\end{tabularx}
\normalsize
\end{center}

\begin{porque}
Se podía escribir todo en un solo archivo enorme o en un notebook gigante. Se descartó porque un archivo de miles de
líneas no se puede revisar ni probar por partes, y un error en un lugar rompe todo lo demás. Con archivos chicos, cada
uno tiene un encabezado que dice: quién lo escribió, qué hace, por qué así, de qué datos sale y \textbf{a qué decisión de
la campaña alimenta}. Si un archivo no podía decir a qué decisión alimenta, no se escribía.
\end{porque}

\section{¿Y los notebooks? ¿Y las pruebas?}
\textbf{Los notebooks} (6) cuentan la historia paso a paso, con texto, tablas y gráficas. No repiten el código: llaman a las
funciones de los cajones. Es como un libro de recetas que dice ``usa la batidora'' en lugar de explicar cómo se fabrica
una batidora.

\begin{center}
\small
\begin{tabular}{ll}
\encabezado \h{Notebook} & \h{Qué cuenta} \\
\codigo{01_planteamiento} & Variables, actores y la desigualdad entre norte y sur \\
\codigo{02_radar} & El índice de presión, los estados y la predicción \\
\codigo{03_pronostico} & La forma del año, los modelos, las tormentas y los escenarios \\
\codigo{05_optimizacion} & El reparto del presupuesto \\
\codigo{06_torre_en_vivo} & La campaña semana a semana \\
\codigo{07_campana} & Las reseñas, las viajeras ideales y la marca \\
\end{tabular}
\normalsize
\end{center}
(No hay notebook 04 porque los escenarios quedaron dentro del 03.)

\textbf{Lo que se entrega a los profesores} es más sencillo: la carpeta \codigo{Entregable} junta esos seis en
\textbf{cuatro notebooks} que funcionan solos, sin archivos \codigo{.py}:
\begin{center}
\small
\begin{tabular}{ll}
\encabezado \h{Notebook de la entrega} & \h{Materias} \\
1. Datos y arquitectura & Grandes volúmenes de datos \\
2. Dónde y cuándo & Planteamiento, minería, aprendizaje de máquina y estocásticos \\
3. Presupuesto y Torre en vivo & Investigación de Operaciones y Spark Streaming \\
4. Campaña & Mercadotecnia digital y minería de texto \\
\end{tabular}
\normalsize
\end{center}
Cada uno trae dentro el código que usa: las celdas que empiezan con \texttt{\%\%modulo} son los ``cajones'' del proyecto,
copiados completos. Así se puede leer todo en un solo lugar y las cifras salen idénticas.

\textbf{Las pruebas} (279) son pequeños programas que revisan solos que todo siga correcto. Hay dos tipos:
\begin{itemize}
\item de \textbf{forma}: que una tabla tenga las columnas que debe, que no haya renglones repetidos;
\item de \textbf{realidad}: que una cifra conocida salga igual. Por ejemplo, que la zona de Tulum tenga 1,031,443
  visitantes en 2025, o que la ocupación de Bacalar en enero de 2024 sea 66.5\,\%.
\end{itemize}
Si alguien vuelve a descargar los datos y una cifra cambia, la prueba falla y avisa. Se corren con un comando:
\texttt{.venv\textbackslash Scripts\textbackslash python -m pytest tests}.

\section{¿Dónde está cada materia (UCA)?}
El proyecto integra las seis materias del semestre. Ninguna está aislada: el resultado de una alimenta a la siguiente.

\begin{center}
\small
\begin{tabularx}{\linewidth}{>{\raggedright\arraybackslash}p{0.2\linewidth}Y>{\raggedright\arraybackslash}p{0.22\linewidth}}
\encabezado \h{Materia} & \h{Qué se hizo, en una frase} & \h{Dónde está} \\
Grandes volúmenes de datos & Se juntaron 8.1 millones de registros oficiales en tres cajas (bronce, plata y oro) y se
limpiaron con Spark; la Torre reproduce 239 semanas con Spark Streaming. & \codigo{base/}, \codigo{envivo/torre.py}
· capítulo~\ref{cap:s-cajas} \\
\filagris Minería de datos & Se encontraron patrones: qué tan concentrado está el turismo, qué meses se llenan, qué
semanas tienen clima raro y qué molesta a los turistas en sus reseñas. & \codigo{radar/indice.py},
\codigo{pronostico/forma.py}, \codigo{campana/texto.py} · capítulo~\ref{cap:s-patrones} \\
Aprendizaje de máquina & Modelos que predicen el estado del mes siguiente y los visitantes de los próximos 12 meses,
comparados contra reglas simples. & \codigo{radar/prediccion.py}, \codigo{pronostico/modelos.py}
· capítulo~\ref{cap:s-predecir} \\
\filagris Procesos estocásticos & Probabilidades: de tormenta cada mes, de que el norte se sature y 10,000 años
simulados para los escenarios. & \codigo{radar/markov.py}, \codigo{pronostico/escenarios.py}
· capítulo~\ref{cap:s-prob} \\
Investigación de Operaciones & Un modelo matemático que reparte el dinero de la campaña de la mejor forma posible
respetando reglas de cuidado. & \codigo{campana/presupuesto.py} · capítulo~\ref{cap:s-dinero} \\
\filagris Mercadotecnia digital & Las viajeras ideales, la marca, los anuncios, los medios y cómo se mide si la campaña
funciona. & \codigo{campana/marca.py} · capítulo~\ref{cap:s-campana} \\
\bottomrule
\end{tabularx}
\normalsize
\end{center}
